
% ==================================================
% 6. NON-FUNCTIONAL REQUIREMENTS
% ==================================================

\section{Non-Functional Requirements}

The following requirements define the main
quality goals and constraints of the
Medication Safety Assistant research prototype.

\begin{longtable}{|
    L{0.13\textwidth}|
    L{0.68\textwidth}|
    L{0.10\textwidth}|
}
\hline

\textbf{ID} &
\textbf{Requirement} &
\textbf{Priority}
\\ \hline

\endfirsthead

\hline
\textbf{ID} &
\textbf{Requirement} &
\textbf{Priority}
\\ \hline

\endhead


% --------------------------------------------------
% NFR-01: PERFORMANCE
% --------------------------------------------------

NFR-01
& \textbf{Performance:}
The system should respond quickly to common
requests. Medication safety checks should
finish within 10 seconds for at least 90\%
of standard test cases under normal conditions.


\PAChange{Performance is measured from
the submission of confirmed medication
information for assessment until the
automated result is displayed.
Tests must define input data, system
conditions, and measurement methods.
HCP response time is excluded.
Longer operations should display
processing status.}

& High
\\ \hline


% --------------------------------------------------
% NFR-02: RELIABILITY AND FAULT TOLERANCE
% --------------------------------------------------

NFR-02
& \textbf{Reliability and Fault Tolerance:}
The system should handle errors and missing
data without crashing.

\PAChange{If OCR, Voice/STT, or NLP fails,
users can retry, correct the input,
or use another available method.}

\PAChange{If some safety checks cannot
be completed, the system must report
their limitations and preserve
available findings. Failed or
unassessed checks must not be
reported as safe.}

\PAChange{Previously detected severe
warnings and related reminder
restrictions must remain recorded
if later checks fail.}
& High
\\ \hline


% --------------------------------------------------
% NFR-03: MEDICATION SAFETY
% --------------------------------------------------

NFR-03
& \textbf{Medication Safety:}

\PAChange{Medication information must be
confirmed by the user before safety assessment.
Confirmation verifies the recorded input,
not medication identity or medical safety.}

\PAChange{Assessment results must distinguish
assessed, partially assessed, unassessed,
and failed states. Missing data or
unsupported checks must not be
presented as proof of safety.}

\PAChange{Detected severe interactions
must trigger immediate warnings.
New reminders for affected medicines
must remain pending the required
HCP review. Unrelated eligible
reminders may proceed.}

\PAChange{Reminder activation requires
a verified schedule from an existing
prescription or recorded HCP instruction.
Proposed changes to prescribed timing
require simulated HCP confirmation
before activation.}

\PAChange{The system must not independently
prescribe, stop, or change medication
treatment. It must never automatically
stop existing reminders because of
a detected interaction.}
& High
\\ \hline


% --------------------------------------------------
% NFR-04: SECURITY AND PRIVACY
% --------------------------------------------------

NFR-04
& \textbf{Security and Privacy:}

\PAChange{The system must protect user
accounts and medication records through
authentication and Role-Based Access
Control (RBAC).}

\PAChange{Patients can access their own
records. HCPs can access only
authorized patient cases and
related conversations.}

\PAChange{Administrators, Dataset Reviewers,
and Technical Support must have access
limited to their assigned tasks.
Technical Support must not receive
unnecessary patient information.}

\PAChange{Sharing patient information
with an HCP requires patient permission.
Access must be limited to the authorized
case, recipient, and shared information.}

\PAChange{The research prototype uses
synthetic or de-identified test data
and does not connect to real
clinical systems.}
& High
\\ \hline


% --------------------------------------------------
% NFR-05: USABILITY AND ACCESSIBILITY
% --------------------------------------------------

NFR-05
& \textbf{Usability and Accessibility:}
The interface should be simple, responsive,
and easy to use. Medication warnings must
be clear, readable, and understandable,
including for users with limited
technical experience.


\PAChange{Warnings must show severity
and limitations using text, not color
alone. Unassessed results, pending reviews,
and failed operations must be clearly
identified.}

\PAChange{Forms and important controls
should have clear labels and readable
text. Confirmation requests and error
messages must explain the required
user action.}

\PAChange{The interface should follow
relevant WCAG 2.2 accessibility
guidelines where feasible for
the research prototype.}

& High
\\ \hline


% --------------------------------------------------
% NFR-06: PLATFORM COMPATIBILITY
% --------------------------------------------------


NFR-06
& \textbf{Platform Compatibility:}

\PAChange{The Web application should
support current versions of major
desktop and mobile browsers, including
Chrome, Edge, Firefox, and Safari,
subject to available test devices.}

Camera, microphone, and notification
features depend on browser support
and user permissions.

\PAChange{Manual input must remain
available when OCR or voice input
cannot be used. Reminder schedules
must remain viewable in the application
when browser notifications are unavailable.}
& High
\\ \hline



% --------------------------------------------------
% NFR-07: DATA INTEGRITY AND HISTORY
% --------------------------------------------------

NFR-07
& \textbf{Data Integrity and History:}
The system should keep medication records,
assessment results, and review history
accurate and consistent. Past results
must be clearly distinguished from
current assessments.

\PAChange{Saved assessments must include
their completion status and time.
Important record changes and review
decisions must retain the responsible
account and relevant timestamps.}

\PAChange{Patient--HCP review conversations,
information requests, replies, and review
outcomes must remain linked to their
authorized review cases.}

\PAChange{New medication information
must not overwrite previous records
without retaining their history.
Earlier safety assessments must not
appear current when their source
medication information has changed.}
& High
\\ \hline


% --------------------------------------------------
% NFR-08: MAINTAINABILITY AND EXTENSIBILITY
% --------------------------------------------------

NFR-08
& \textbf{Maintainability and Extensibility:}
The system should use a modular design
to support updates, new data sources,
and future medication safety features
without major changes to existing modules.

\PAChange{Input processing, medication
identification, safety assessment,
HCP review, and reminder management
should have clearly separated
responsibilities.}
& Medium
\\ \hline


% --------------------------------------------------
% NFR-09: DOCUMENTATION
% --------------------------------------------------


NFR-09
& \textbf{Documentation:}

\PAChange{The project should provide
installation and setup instructions,
a user guide for main features,
and technical documentation
for developers.}

\PAChange{Documentation should explain
supported safety checks, data coverage,
missing-data handling, and the simulated
HCP review workflow.}

\PAChange{It must clearly state that
the prototype is for research and
educational use, not independent
clinical decision-making.}
& Medium
\\ \hline



% --------------------------------------------------
% NFR-10: DATA SOURCE TRACEABILITY
% --------------------------------------------------

\PAChange{NFR-10}
& \PAChange{\textbf{Data Source Traceability:}
Medication data and ingredient mappings
used for safety assessments should retain
available source references and
dataset review status.}

\PAChange{Safety assessments must distinguish
supported findings from missing or
unverified source information.
Dataset review decisions must not
be presented as independent
clinical validation.}

\PAChange{Data use must follow applicable
source terms and permissions.}
& \PAChange{High}
\\ \hline

\end{longtable}
